\section{Introduction and problem statement}
\label{sec:intro}

\subsection{Context: an e-commerce customer service centre}
Most contacts to an online shop's customer service are caused by orders: a parcel is late (``where is my
order?''), a customer wants to cancel or get a refund, an item arrives damaged or wrong, or an invoice or
address must be changed. These causes are largely \emph{predictable}: the number of orders, the promotion
calendar and the carriers' delays are known days in advance. Yet many centres still plan from last year's
averages and react to peaks when they happen. This status quo has four costs:
\begin{itemize}
  \item \textbf{Peaks hurt customers}: long waits, abandoned chats, repeat contacts and bad reviews.
  \item \textbf{Quiet periods waste money}: staffing is the largest cost of a contact centre, and paid advisors
        sit idle.
  \item \textbf{Decisions are tested on real customers}: there is no safe place to ask ``what if a promotion
        doubles the tracking questions?'' or ``what if 20\% of the advisors are sick?''.
  \item \textbf{AI assistants are adopted without measurement}: chatbots answer customers, but few companies
        check whether the answers follow the company policy, or whether an automatic grader agrees with humans.
\end{itemize}

\subsection{Our answer: a digital twin with four agents}
CallVerse is a \emph{digital twin} of such a centre: a virtual copy that behaves like the real one, so that
a decision can be tested in seconds before customers feel it. The twin is driven by real e-commerce data
(the Olist marketplace) and is used by four cooperating \emph{agents}, programs that observe a situation and
act on it:
\begin{itemize}
  \item the \textbf{Client agent} plays customers, each with a real order, a problem and a personality;
  \item the \textbf{Workforce agent} recognises each customer's problem, sends the chat to the right advisor,
        and plans the advisors' shifts;
  \item the \textbf{Advisor agent} answers, using the company's policy documents and business tools (order
        lookup, refund eligibility, tickets);
  \item the \textbf{Quality agent} grades every conversation on a five-criterion scorecard.
\end{itemize}

\subsection{Business and data science objectives}
We framed the project as four \emph{business objectives} (BO), each measured by business KPIs, and four
\emph{data science objectives} (DSO), each measured by a model metric (Table~\ref{tab:bodso}). A KPI tells
whether the centre is doing better; a DSO metric tells whether the AI component works. Three rules apply to
every objective: (1) the main idea must beat a \emph{simple baseline}; (2) a \emph{backup} covers the whole
objective with less effort; (3) a \emph{switch rule}, fixed before the experiments, decides whether the main
idea or the backup goes to production. DSO4 is the foundation: DSO1 and DSO3 are trained and tested inside
its simulation.

\begin{table}[h]
\centering\small
\caption{From business objectives to data science objectives. ``Rule'' is the switch rule decided in advance.}
\label{tab:bodso}
\begin{tabularx}{\textwidth}{@{}p{2.5cm}YYYp{3.1cm}@{}}
\toprule
\textbf{Business objective} & \textbf{DSO and main idea} & \textbf{Baseline} & \textbf{Backup} & \textbf{Rule} \\
\midrule
BO1 Customers wait less & DSO1: real-time routing learned by reinforcement learning (MaskablePPO) &
  First come, first served (FIFO) with skill matching & Priority score with tuned weights &
  $\geq$10\% less waiting in the twin's peak weeks \\
BO2 Problems solved correctly, higher satisfaction & DSO2: advisor with retrieval-augmented generation and
  tools, plus an LLM judge & Keyword FAQ search; templates & Intent classifier + approved templates +
  checklist reviewer & $\leq$10\% wrong answers; judge--human kappa $\geq$ 0.4 \\
BO3 Lower cost, same service & DSO3: forecast + constraint-programming shift planning (CP-SAT) &
  Fixed schedule sized for the peak & Erlang C + shift templates + overtime rule &
  $\geq$5\% cheaper while meeting the service target and the labour rules \\
BO4 Prepare for peaks by testing first & DSO4: contact forecast per reason + discrete-event simulation &
  Same day last week; Erlang C queueing formula & Statistical forecast + queueing model &
  Beat the naive forecast; the twin passes its fidelity checks \\
\bottomrule
\end{tabularx}
\end{table}

\paragraph{What success means.} The centre's service target is the industry's ``80/20'' rule: 80\% of chats
answered within 20 seconds. A good plan meets it at the lowest labour cost and respects labour law (at most
5 days and 40 hours a week, at least 11 hours of rest between two shifts). A good answer applies the company
policy (for example: a parcel more than 7 days late triggers a carrier investigation) and leaves the customer
satisfied. Every claimed gain must be measured against its baseline on the same simulated customers, with a
95\% confidence interval.

\subsection{Contributions and structure of the report}
We deliver (1) a reproducible data pipeline from raw Olist and Bitext files to a relational store and a
knowledge base; (2) a forecast and a validated simulation engine; (3) four agents combining machine learning,
reinforcement learning, optimisation and language models; (4) an evaluation framework with paired
experiments, a scorecard and live monitoring of every AI component; and (5) a deployed application.
Section~\ref{sec:data} presents the data, Section~\ref{sec:method} the methods and architecture,
Section~\ref{sec:experiments} the experimental protocol and Section~\ref{sec:results} the results.
Section~\ref{sec:analysis} discusses limitations, Section~\ref{sec:deployment} deployment and
reproducibility, and Section~\ref{sec:aiuse} declares our use of generative AI tools.
